\bookchapter{נספח ד': איפה למצוא משאבים אמינים?}

\booksection{תמיכה רשמית היא הבחירה הראשונה:}

\begin{itemize}
  \item \textbf{משתמשי אנדרואיד}: באתרים הרשמיים של מותגי הטלפונים יש מדור "שירות ותמיכה".
  \item \textbf{כלים כלליים}: Apple – iTunes, Huawei – Huawei Phone Assistant, Xiaomi – Mi Mover.
  \item \textbf{אתרים ויוצרי תוכן שימושיים:}חיפוש טיפים
\end{itemize}

\booksection{אתרים ויוצרי תוכן שימושיים:}

\begin{itemize}
  \item \textbf{צפייה בביקורות}: צפו במספר סרטוני ביקורת של יוצרי תוכן שונים.
  \item \textbf{לימוד צילום}: עקבו אחר יוצרי תוכן שמלמדים צילום בטלפון.
  \item \textbf{רשימת תיוג אבטחה במשפט אחד:}הגדרתם סיסמת נעילה?
\end{itemize}

\booksection{רשימת תיוג אבטחה במשפט אחד:}

\begin{enumerate}
  \item הפעלתם אימות דו-שלבי לחשבונות החשובים?
  \item מגבים באופן קבוע תמונות וצ'אטים של WeChat?
  \item זכרו את מנטרת ההגנה מפני הונאות: אל תאמינו בקלות, אל תחשפו מידע, אל תעבירו כסף, אל תהיו חמדנים, אל תמהרו.
  \item נספח ה': תודות והערות
  \item תודה שהתמדתם בקריאת "שיעור הטלפון" הזה! בתהליך הכתיבה, התייעצנו במדריכים רשמיים של יצרנים, ספרי הדרכה למשתמש ומאמרי טכנולוגיה פומביים כדי להבטיח את דיוק המידע.
\end{enumerate}
